\documentclass{ximera}
% xmPreamble.tex en xmPrintstyle.sty worden automatisch geladen door ximera.cls (v2.7.8);
% NIET nog eens \input{./preamble.tex} of \addPrintStyle{.} doen (geeft 'already defined'-fouten).

\makeatletter % eigen img-map toevoegen aan het zoekpad (werkt standalone én vanuit de xourse)
\g@addto@macro\Ginput@path{{img/}{UW4202/img/}}
\makeatother

\begin{document}
\author{Alexander Holvoet}
\xmtitle{Convexe grondvlakken: trapezia}{}

\begin{instructorNotes}
    Dit stuk is in de bron niet volledig als lesactiviteit uitgeschreven, maar het is wel
    de bedoeling om dit op een onderzoekende manier aan te brengen, zoals in de vorige
    paragrafen. Dit kan in groepjes of via een onderwijsleergesprek, omdat het wat
    moeilijker wordt. Het interessante van deze situatie is niet zozeer het eindresultaat,
    maar vooral de zoektocht ernaartoe.
\end{instructorNotes}

We bekijken enkele trapezia. Een trapezium is een vierhoek met twee evenwijdige zijden.
We vertrekken van een driehoek: teken de bissectrices en de ingeschreven cirkel. Doe dit
in viervoud. Knip nu van drie van die driehoeken een stuk af volgens een lijn evenwijdig
met de basis: boven de ingeschreven cirkel (situatie 2), rakend aan de ingeschreven cirkel
(situatie 3) en dwars door de ingeschreven cirkel (situatie 4). Teken ten slotte ook de
bissectrices van de twee nieuw verschenen hoeken van elk trapezium.

% TODO (afbeelding ontbreekt): \afbeelding[15cm]{Ltrapezia.png} (een driehoek en drie
% soorten trapezia). Lijntekening niet aanwezig in img/.

Voor je het experiment uitvoert, kun je proberen te voorspellen welke zandhoop en welk
bovenaanzicht je zult krijgen. Het loont echt de moeite om dit zelf verder uit te zoeken.
Als je zand giet op de volledige driehoek (situatie 1), krijg je uiteraard een piramide.
We bekijken situatie 2 en 4 in detail; situatie 3 laten we aan de lezer.

\textbf{Situatie 2.} In de tweede situatie zijn de bissectrices niet concurrent. Er zijn
twee snijpunten van telkens twee bissectrices, de punten \(O\) en \(O'\). Ze liggen op de
bissectrice van de afgesneden hoek van de driehoek en zijn de projecties van \(S\) en
\(S'\), de eindpunten van de kamlijn.

In de 3D-tekening zien we twee rechthoekige driehoeken \(HSO\) en \(H'S'O'\) die dezelfde
rusthoek hebben en dus gelijkvormig zijn. We zien dat \(|OS| = |OH|\cdot \tan\theta\) en
\(|O'S'| = |O'H'|\cdot \tan\theta\), maar \(|OH| > |O'H'|\), dus \(|OS| > |O'S'|\). Het
lijnstuk \([SS']\) loopt dus niet horizontaal. In het bovenaanzicht zie je dit aan het feit
dat \(O'\) dichter bij de rand ligt dan \(O\). Verifieer het met het experiment. Sterker
nog:
\[
    |OS| - |O'S'| = (|OH| - |O'H'|)\cdot \tan\theta .
\]
Als we de kniplijn helemaal zouden opschuiven tot in het hoekpunt, dan zou \(|O'H'| = 0\)
worden, dus de helling van \([SS']\) zou dezelfde zijn als die van \([SA]\). Als we
\([SS']\) verlengen, zou die door de top \(A\) van de afgeknipte driehoek lopen.

\begin{image}[7cm]
    \includegraphics[width=7cm,keepaspectratio]{lsituatie2.jpg}
\end{image}
Situatie 2: het afgeknipte trapezium met de kamlijn.

% TODO (afbeeldingen ontbreken): in de bron stonden hier de lijntekeningen
% Ltrap2.png (projectie situatie 2) en Ltrap23D.png (3D-versie). Niet aanwezig in img/.

\textbf{Situatie 4.} Hier hebben we ook de vier bissectrices van de hoeken van het
trapezium getekend. Nu snijdt telkens \'e\'en van de oude bissectrices met \'e\'en van de
nieuwe bissectrices. Dit geeft weer twee punten \(O\) en \(O'\). Het lijnstuk \([OO']\) dat
deze punten verbindt, is evenwijdig met de twee evenwijdige zijden van het trapezium.
\(O\) en \(O'\) zijn de projecties op het grondvlak van de toppen \(S\) en \(S'\).

In de 3D-voorstelling zijn er weer twee driehoeken \(HSO\) en \(H'S'O'\) die dezelfde
rusthoek met het grondvlak hebben. We vinden \(|OS| = |OH|\cdot \tan\theta\) en
\(|O'S'| = |O'H'|\cdot \tan\theta\), maar nu is \(|OH| = |O'H'|\) en dus \(|OS| = |O'S'|\).
Bijgevolg is \([SS'] \parallel [HH'] \parallel [KK']\): de kamlijn is evenwijdig met de
evenwijdige zijden van het trapezium. Verifieer weer met het experiment.

\begin{image}[7cm]
    \includegraphics[width=7cm,keepaspectratio]{lsituatie4.jpg}
\end{image}
Situatie 4: de kamlijn is evenwijdig met de evenwijdige zijden van het trapezium.

% TODO (afbeeldingen ontbreken): in de bron stonden hier Ltrap4.png (projectie situatie 4)
% en Ltrap43D.png (3D-tekening). Niet aanwezig in img/.

\end{document}
